% Copyright 2026 Sreeram Anil
% SPDX-License-Identifier: CC-BY-4.0
\chapter{Introduction}
\label{ch:intro}

\section{Why a battery estimator benchmark for electric aircraft}

An electric aircraft cannot be flown on the quantities its sensors measure. The pilot and
the flight-control software need the state of charge (the fraction of the usable charge
that remains), the state of health (the fraction of the original capacity that the cell can
still deliver), the remaining energy and the power that can be drawn for the next
manoeuvre. None of these is observable directly; all of them are inferred from terminal
voltage, current and temperature through a model of the electrochemical cell, and the
quality of that inference decides how much of the battery must be held back as reserve.
For a vertical take-off vehicle whose landing hover draws four to five times the cruise
current \citep{yang2021evtol,fredericks2018vtol}, a two-percent error in the estimated state
of charge at the end of a mission is not an accounting detail: it is the difference between
a reserve that exists and one that does not. Certification standards for rechargeable
lithium batteries in aircraft \citep{rtca2017do311a} consequently expect the state estimate
to be accurate, bounded, and predictable under sensor faults and cell ageing.

The estimation problem itself is a classic of control engineering. The cell is a nonlinear
dynamic system with a slowly varying parameter set, the measurement is scalar and only
weakly informative on the flat parts of the open-circuit-voltage curve, the sensors carry
biases that integrate into unbounded errors, and the whole system operates over a
temperature range in which every impedance parameter changes by an order of magnitude.
Practically every estimator ever proposed in the control literature --- Kalman filters in
all their variants, deterministic observers, particle methods, moving-horizon estimation,
learning-based hybrids --- has been applied to it, and a battery-management engineer is
expected to know which of them to use, why, and at what computational cost on a
microcontroller that must also run the rest of the BMS.

This report and the accompanying software answer that question empirically. Rather than
describing the estimators, it implements all of them, in one code base and against one
interface, and runs them on the same missions, faults and cells, so that every comparison
in the closing chapters is a controlled experiment rather than a literature summary.

\section{What was built}

The software is \estkit, a header-only C++17 library organised around a small model
concept (\cref{ch:library}). A model is any class that provides a state transition
$\vx_{k+1}=\bm f(\vx_k,\vu_k)$, a measurement $\vy_k=\bm h(\vx_k,\vu_k)$ and the two noise
covariances; Jacobians, constraints and tunable parameters are optional and are supplied
numerically when absent. Every filter and observer in the library is a template over such a
model, so the same \code{Ekf<M>}, \code{Ukf<M>} or \code{Mhe<M,N>} that estimates the state
of charge of a cell estimates the position of a drone or the attitude of an airframe when
handed the corresponding model. The battery model, the plants that generate the benchmark
data, the fault scenarios and the harness are separate from the estimators and are the
only battery-specific parts of the code.

The library contains, as of this report, 70 cell-level estimators in ten families, seven
pack-level estimators, nine attitude estimators and the tools around them; every one of
them is implemented from its primary reference, cited in the header of its source file and
in its chapter. The design constraints are those of flight software: no dynamic memory, no
exceptions, no recursion of unbounded depth, fixed-size linear algebra with compile-time
dimensions, deterministic execution, a single \code{Real} type that is \code{double} on the
benchmark machine and \code{float} on a microcontroller, and warning-free compilation with
\code{-Wall -Wextra -Wpedantic -Werror} on GCC and Clang.

The benchmark generates its data from two truth plants that are deliberately richer than
the estimators' model (\cref{ch:ecm,ch:spm}): an enhanced equivalent-circuit cell with a
lumped thermal state, Arrhenius impedance scaling, one-state and Preisach hysteresis and
cycle/calendar ageing, and an electrochemical single-particle model built from the
published electrode data of the LG~M50 cell \citep{chen2020parameterization} and validated
against PyBaMM \citep{sulzer2021pybamm}. Four mission profiles, twelve fault and operating
scenarios and three noise realisations produce the cell-level data set (\cref{ch:datasets});
a twelve-cell series module, a 120-mission ageing study and a synthetic flight with a
nine-axis inertial unit produce the pack, state-of-health and attitude data sets.

\section{How to read this report}

\Cref{ch:prelim} fixes the notation and collects the estimation theory that the estimator
chapters build on. \Cref{ch:ecm,ch:spm} derive the two cell models, \cref{ch:datasets} the
missions and scenarios, \cref{ch:metrics} the metrics and the validation against external
reference libraries, and \cref{ch:library} the software architecture and the embedded
considerations.

Parts~II to~XIII contain one chapter per estimator. Each chapter follows the same template:
the motivation and the aerospace or BMS context; the problem statement with its assumptions;
the derivation with every symbol defined where it first appears; an algorithm box; the
instantiation on the battery model; implementation notes covering numerical safeguards,
memory, cost and single-precision behaviour; and the benchmark results with the
auto-generated table and figures for that estimator and a discussion of where and why it
succeeds or fails. The chapters are self-contained and can be read in any order, but within
a family the later chapters assume the earlier ones.

Part~XIV synthesises the results: the cell-level ranking on the equivalent-circuit plant,
the behaviour under the structured model error of the electrochemical plant, the numerical
robustness in single precision, the pack, attitude and multi-flight results, and finally
recommendations for an electric-aircraft BMS and a discussion of what the benchmark does not
cover.

\section{Scope and honesty}

Three limitations are stated here and repeated where they matter. First, the data are
simulated. The plants are built from published electrode data and literature model
structures, and the impedance, thermal, ageing and sensor parameters are representative of
a high-power NMC cell rather than measured on one; the parameter tables in \cref{app:params}
label the provenance of every value. The absolute numbers in the results tables therefore
describe this benchmark, not a particular product; the relative behaviour of the estimators
--- which of them is hurt by which fault, and by how much --- is the transferable content.
Second, each estimator was implemented and tuned by an engineer with the primary reference
in hand and with literature-standard settings; no estimator was tuned on the test data
beyond what a practitioner would do, but neither was any of them tuned by its inventor.
Third, the execution times were measured on a single x86-64 core; the projection to a
Cortex-M class processor in \cref{ch:metrics} is an estimate with stated assumptions, not a
measurement.

Everything in this report is reproducible: \cref{app:build} lists the commands that rebuild
the library, re-run every benchmark, regenerate every figure and table from the results, and
re-run the cross-checks against FilterPy, PyBaMM and \texttt{ahrs}.
